\chapter{1985年上海卷（文）}

\section{填空题}

\begin{problem}

不等式 \(\left|x+2\right|<1\) 的解集是\fillinblank{}.

\begin{answer}
\(\left(-3,-1\right)\)
\end{answer}

\end{problem}

\begin{solution}
由绝对值不等式的意义，\(\left|x+2\right|<1\) 等价于 \(-1<x+2<1\).
不等式两边同时减去 \(2\)，得 \(-3<x<-1\)，所以解集为 \(\left(-3,-1\right)\).
\end{solution}

\begin{problem}

函数 \(y=\sqrt[8]{x}-1\) 的反函数是\fillinblank{}.

\begin{answer}
\(y=(x+1)^8\)，\(x\geqslant -1\)
\end{answer}

\end{problem}

\begin{solution}
原函数的定义域为 \(x\geqslant0\)，值域为 \(y\geqslant-1\). 由 \(y=\sqrt[8]{x}-1\) 得 \(y+1=\sqrt[8]{x}\)，两边取八次方，得 \(x=(y+1)^8\). 交换 \(x\)，\(y\) 的位置，得到反函数 \(y=(x+1)^8\)，其定义域是原函数的值域，即 \(x\geqslant-1\).
\end{solution}

\begin{problem}

函数 \(y=\tan 2x\) 的最小正周期是\fillinblank{}.

\begin{answer}
\(\dfrac{\pi}{2}\)
\end{answer}

\end{problem}

\begin{solution}
正切函数 \(\tan u\) 的最小正周期为 \(\pi\). 设 \(y=\tan 2x\) 的正周期为 \(T\)，则 \(2T=\pi\)，所以 \(T=\dfrac{\pi}{2}\). 因而最小正周期为 \(\dfrac{\pi}{2}\).
\end{solution}

\begin{problem}

设等差数列 \(\{a_n\}\) 中，\(a_1=2\)，\(a_2=5\)，则 \(a_{10}=\)\fillinblank{}.

\begin{answer}
\(29\)
\end{answer}

\end{problem}

\begin{solution}
等差数列的公差为 \(d=a_2-a_1=5-2=3\). 因此 \(a_{10}=a_1+9d=2+9\times3=29\).
\end{solution}

\begin{problem}

方程 \(9^x=3^{x+1}\) 的解是\fillinblank{}.

\begin{answer}
\(x=1\)
\end{answer}

\end{problem}

\begin{solution}
将 \(9\) 化为以 \(3\) 为底，得 \(9^x=3^{2x}\). 由于底数 \(3>0\) 且 \(3\neq1\)，指数相等，即 \(2x=x+1\)，解得 \(x=1\).
\end{solution}

\begin{problem}

用不等号“>”或“<”连接：\(\left(\frac{1}{2}\right)^{\frac{1}{7}}\)\fillinblank{}\(\left(\frac{1}{2}\right)^{\frac{1}{6}}\).

\begin{answer}
\(>\)
\end{answer}

\end{problem}

\begin{solution}
因为 \(0<\dfrac{1}{2}<1\)，函数 \(y=\left(\dfrac{1}{2}\right)^t\) 在实数范围内单调递减. 又因为 \(\dfrac{1}{7}<\dfrac{1}{6}\)，所以 \(\left(\frac{1}{2}\right)^{\frac{1}{7}}>\left(\frac{1}{2}\right)^{\frac{1}{6}}\).
\end{solution}

\begin{problem}

如果一个圆的圆心在点 \((2,4)\)，并且经过点 \((0,3)\)，则这个圆的方程是\fillinblank{}.

\begin{answer}
\(\left(x-2\right)^2+\left(y-4\right)^2=5\)
\end{answer}

\end{problem}

\begin{solution}
圆心为 \(C(2,4)\)，半径的平方为
\[
r^2=(0-2)^2+(3-4)^2=4+1=5
\]
所以圆的方程为 \(\left(x-2\right)^2+\left(y-4\right)^2=5\).
\end{solution}

\begin{problem}

直线 \(y=-x+3\) 的倾斜角的大小是\fillinblank{}.

\begin{answer}
\(\dfrac{3\pi}{4}\)
\end{answer}

\end{problem}

\begin{solution}
直线的斜率为 \(k=-1\). 设倾斜角为 \(\theta\)，则 \(0\leqslant\theta<\pi\) 且 \(\tan\theta=k=-1\). 在这个范围内，满足条件的角为 \(\theta=\dfrac{3\pi}{4}\).
\end{solution}

\begin{problem}

设 \(\lg 2=a\)，则 \(\log_2 25=\)\fillinblank{}.

\begin{answer}
\(\dfrac{2(1-a)}{a}\)
\end{answer}

\end{problem}

\begin{solution}
由换底公式，\(\log_2 25=\frac{\lg25}{\lg2}=\frac{2\lg5}{a}\). 又因为 \(\lg5=\lg\dfrac{10}{2}=1-\lg2=1-a\)，所以 \(\log_2 25=\frac{2(1-a)}{a}\).
\end{solution}

\begin{problem}

已知 \(O\) 为直角坐标系的原点，\(A\) 点的坐标为 \((1,3)\)，\(B\) 点的坐标为 \((0,3)\)，若把 \(\triangle OAB\) 绕 \(y\) 轴旋转一周，则所得旋转体的体积是\fillinblank{}.

\begin{answer}
\(\pi\)
\end{answer}

\end{problem}

\begin{solution}
线段 \(OA\) 的方程为 \(x=\dfrac{y}{3}\)，其中 \(0\leqslant y\leqslant3\). 三角形 \(OAB\) 在高度 \(y\) 处的截面是从 \(y\) 轴到直线 \(OA\) 的线段，绕 \(y\) 轴旋转后得到半径 \(\dfrac{y}{3}\) 的圆盘. 因此旋转体的体积为
\[
V=\int_0^3\pi\left(\frac{y}{3}\right)^2\,\mathrm{d}y
=\frac{\pi}{9}\left[\frac{y^3}{3}\right]_0^3
=\pi
\]
\end{solution}

\section{单选题}

\begin{problem}

函数 \(y=-\sin 2x+3\) 的最大值是（\quad）

\choices
    {\(2\)}
    {\(3\)}
    {\(4\)}
    {\(5\)}

\begin{answer}
C
\end{answer}

\end{problem}

\begin{solution}
因为 \(-1\leqslant\sin2x\leqslant1\)，所以 \(-1\leqslant-\sin2x\leqslant1\). 从而 \(2\leqslant-\sin2x+3\leqslant4\).
当 \(\sin2x=-1\) 时取到上界，因此最大值为 \(4\)，选 C.
\end{solution}

\begin{problem}

设复数 \(z_1=1+\i\)，\(z_2=1-3\i\)，则复数 \(z=z_1\cdot z_2\) 在复平面内所表示的点位于（\quad）

\choices
    {第一象限}
    {第二象限}
    {第三象限}
    {第四象限}

\begin{answer}
D
\end{answer}

\end{problem}

\begin{solution}
计算得
\[
z=(1+\i)(1-3\i)=1-3\i+\i-3\i^2=4-2\i
\]
复数对应的点为 \((4,-2)\)，实部为正、虚部为负，所以点位于第四象限，选 D.
\end{solution}

\begin{problem}

已知函数 \(f(x)=\lg x^2\) 的定义域为 \(F\)，函数 \(g(x)=2\lg x\) 的定义域为 \(G\)，那么（\quad）

\choices
    {\(F \cap G = \varnothing\)}
    {\(F = G\)}
    {\(F \subset G\)}
    {\(G \subset F\)}

\begin{answer}
D
\end{answer}

\end{problem}

\begin{solution}
函数 \(f(x)=\lg x^2\) 有意义的条件是 \(x^2>0\)，即 \(x\neq0\)，所以 \(F=\{x\mid x\neq0\}\). 函数 \(g(x)=2\lg x\) 有意义的条件是 \(x>0\)，所以 \(G=(0,+\infty)\). 显然 \(G\subset F\)，选 D.
\end{solution}

\begin{problem}

若一个四棱锥的底面是凸四边形，它的顶点到底面各边的距离都相等，则这个棱锥的底面四边形（\quad）

\choices
    {必为正方形}
    {必有内切圆}
    {必有外接圆}
    {必既有内切圆又有外接圆}

\begin{answer}
B
\end{answer}

\end{problem}

\begin{solution}
设棱锥顶点为 \(P\)，底面所在平面为 \(\alpha\)，点 \(P\) 在 \(\alpha\) 上的射影为 \(H\). 对底面任一边所在直线 \(l\)，由 \(PH\perp\alpha\) 可得
\[
\operatorname{dist}(P,l)^2=PH^2+\operatorname{dist}(H,l)^2
\]
题设中 \(\operatorname{dist}(P,l)\) 对四条边相等，且 \(PH\) 是定值，所以点 \(H\) 到底面四条边所在直线的距离相等. 由于底面是凸四边形，点 \(H\) 是其内心，取以 \(H\) 为圆心、以该公共距离为半径的圆，它同时与四条边相切. 因此底面四边形必有内切圆，选 B.
\end{solution}

\begin{problem}

\(a>0\) 且 \(b>0\)，是 \(ab>0\) 的（\quad）

\choices
    {充分但不必要的条件}
    {必要但不充分的条件}
    {充分而且必要的条件}
    {既不充分又不必要的条件}

\begin{answer}
A
\end{answer}

\end{problem}

\begin{solution}
若 \(a>0\) 且 \(b>0\)，则 \(ab>0\)，所以该条件是充分条件. 但 \(ab>0\) 还可能由 \(a<0\) 且 \(b<0\) 得到，因此 \(a>0\) 且 \(b>0\) 不是必要条件. 所以它是充分但不必要的条件，选 A.
\end{solution}

\section{解答题}

\begin{problem}

在平面直角坐标系内，画出下列方程的曲线（只要求画出图形）：

\begin{enumerate}
\item \(\frac{x}{3} - \frac{y}{2} = -1\)；
\item \(\frac{x^2}{9} - \frac{y^2}{4} = -1\).
\end{enumerate}
\end{problem}

\begin{solution}
\begin{enumerate}
\item 将方程化为斜截式，得 \(y=\dfrac{2}{3}x+2\). 这是一条斜率为 \(\dfrac{2}{3}\) 的直线，且它与坐标轴分别交于 \((-3,0)\) 和 \((0,2)\). 在坐标系中标出这两点并连接，即得所求直线.
\item 将方程移项，得
\[
\frac{y^2}{4}-\frac{x^2}{9}=1
\]
这是以原点为中心、实轴在 \(y\) 轴上的双曲线，顶点为 \((0,2)\) 和 \((0,-2)\)，渐近线为 \(y=\dfrac{2}{3}x\) 与 \(y=-\dfrac{2}{3}x\). 据此即可画出所求双曲线.
\end{enumerate}
\end{solution}

\begin{problem}

\begin{enumerate}
\item 求 \(\left(2+x^2\right)^8\) 的展开式中 \(x^{10}\) 的系数.

\item 从六个数字 \(1\)，\(2\)，\(3\)，\(5\)，\(7\)，\(9\) 中任取四个不同的数字，有多少种取法？由这六个数字可以组成多少个没有重复数字的四位偶数？

\item 在一列球中，第 \(1\) 个球的半径为 \(1\)，第 \(2\) 个球的直径是第 \(1\) 个球的半径，以后每一个球的直径都是前一个球的半径，这样无限继续下去，求所有这些球的表面积的和.
\end{enumerate}
\end{problem}

\begin{solution}
\begin{enumerate}
\item 二项式展开的通项中，若从 \(8\) 个因子中选出 \(k\) 个 \(x^2\)，则该项为 \(\mathrm C_8^k2^{8-k}x^{2k}\). 要得到 \(x^{10}\)，必须有 \(2k=10\)，即 \(k=5\). 因此所求系数为 \(\mathrm C_8^5 2^{8-5}=56\times8=448\).

\item 从六个数字中任取四个不同的数字，只需从六个元素中选取四个，取法数为 \(\mathrm C_6^4=\frac{6\times5\times4\times3}{4\times3\times2\times1}=15\). 在所给六个数字中，只有数字 \(2\) 是偶数，所以四位偶数的个位必须是 \(2\). 个位确定后，千位、百位、十位从其余五个数字中取出三个并按顺序排列，组成数的种数为 \(\mathrm A_5^3=5\times4\times3=60\). 所以两问的答案分别为 \(15\) 种和 \(60\) 个.

\item 记第 \(n\) 个球的半径为 \(r_n\). 已知 \(r_1=1\)，且从第二个球开始，\(2r_n=r_{n-1}\)，所以 \(r_n=\frac{1}{2^{n-1}}\). 第 \(n\) 个球的表面积为 \(S_n=4\pi r_n^2=4\pi\left(\frac{1}{2^{n-1}}\right)^2=4\pi\left(\frac{1}{4}\right)^{n-1}\). 于是所有球的表面积之和是首项为 \(4\pi\)、公比为 \(\dfrac14\) 的无穷等比数列之和，故 \(S=\frac{4\pi}{1-\frac14}=\frac{16\pi}{3}\).
\end{enumerate}
\end{solution}

\begin{problem}

如图，直线 \(a\parallel b\)，点 \(P\) 在 \(a\)，\(b\) 所确定的平面外，\(PA\perp a\)，且交 \(a\) 于 \(A\) 点；\(AB\perp b\)，且交 \(b\) 于 \(B\) 点. 求证：\(PB\perp b\).
\end{problem}

\begin{solution}
设过 \(P\)，\(A\)，\(B\) 的平面为 \(\beta\). 因为 \(AB\perp b\)，且 \(a\parallel b\)，所以 \(AB\perp a\). 又已知 \(PA\perp a\)，而 \(PA\) 与 \(AB\) 是平面 \(\beta\) 内相交的两条直线，故 \(a\perp\beta\).

由 \(a\parallel b\)，得 \(b\perp\beta\). 直线 \(PB\) 在平面 \(\beta\) 内，且与 \(b\) 交于 \(B\)，所以 \(PB\perp b\).
\end{solution}

\begin{problem}

\begin{enumerate}
\item 证明万能公式：\(\sin\alpha=\frac{2\tan\frac{\alpha}{2}}{1+\tan^2\frac{\alpha}{2}}\).

\item 已知：\(\cos\alpha=-\frac{1}{3}\)，\(\pi<\alpha<\frac{3\pi}{2}\)，求：\(\cos\left(\frac{\alpha}{2}+\frac{5\pi}{6}\right)\).
\end{enumerate}
\end{problem}

\begin{solution}
\begin{enumerate}
\item 在 \(\tan\dfrac{\alpha}{2}\) 有意义的条件下，由正弦的二倍角公式，\(\sin\alpha=2\sin\frac{\alpha}{2}\cos\frac{\alpha}{2}\). 将右端的分子、分母同除以 \(\cos^2\dfrac{\alpha}{2}\)，得到 \(\sin\alpha=\frac{2\sin\frac{\alpha}{2}\cos\frac{\alpha}{2}}{\cos^2\frac{\alpha}{2}}=\frac{2\tan\frac{\alpha}{2}}{\sec^2\frac{\alpha}{2}}\). 又因为 \(\sec^2 t=1+\tan^2t\)，所以 \(\sin\alpha=\frac{2\tan\frac{\alpha}{2}}{1+\tan^2\frac{\alpha}{2}}\). 公式得证.

\item 由 \(\pi<\alpha<\dfrac{3\pi}{2}\)，可知 \(\dfrac{\pi}{2}<\dfrac{\alpha}{2}<\dfrac{3\pi}{4}\)，所以 \(\dfrac{\alpha}{2}\) 位于第二象限，正弦为正、余弦为负. 由半角公式，\(\cos\frac{\alpha}{2}=-\sqrt{\frac{1+\cos\alpha}{2}}=-\sqrt{\frac{1-\frac13}{2}}=-\frac{1}{\sqrt3}\)，\(\sin\frac{\alpha}{2}=\sqrt{\frac{1-\cos\alpha}{2}}=\sqrt{\frac{1+\frac13}{2}}=\sqrt{\frac23}\). 因此
\[
\cos\left(\frac{\alpha}{2}+\frac{5\pi}{6}\right)
=\left(-\frac{1}{\sqrt3}\right)\left(-\frac{\sqrt3}{2}\right)-\sqrt{\frac23}\cdot\frac12
=\frac12-\frac{\sqrt6}{6}=\frac{3-\sqrt6}{6}
\]
\end{enumerate}
\end{solution}

\begin{problem}

设 \(\triangle ABC\) 的两个内角 \(A\)，\(B\) 所对的边分别为 \(a\)，\(b\). 若复数 \(z_1\) 的模为 \(a\)，幅角为 \(B\)；复数 \(z_2\) 的模为 \(b\)，幅角为 \(-A\). 求证：\(z_1+z_2\) 为实数.
\end{problem}

\begin{solution}
由复数的三角形式，\(z_1=a\left(\cos B+\i\sin B\right)\)，\(z_2=b\left(\cos A-\i\sin A\right)\). 所以 \(z_1+z_2=\left(a\cos B+b\cos A\right)+\i\left(a\sin B-b\sin A\right)\).
在 \(\triangle ABC\) 中，由正弦定理，存在同一个正数 \(k\)，使得 \(a=k\sin A\)，\(b=k\sin B\). 因此 \(a\sin B-b\sin A=k\sin A\sin B-k\sin B\sin A=0\).
于是 \(z_1+z_2=a\cos B+b\cos A\)，其虚部为零，所以 \(z_1+z_2\) 为实数.
\end{solution}

\begin{problem}

如图，由抛物线 \(y=x^2+2\) 上的点 \(M(x_0,y_0)\) 向直线 \(y=\frac{1}{2}x\) 作垂线，垂足为 \(N\)，延长 \(MN\) 至 \(P\) 点，使点 \(N\) 分线段 \(MP\) 成定比 \(\lambda=\frac{MN}{NP}=4\).

\begin{enumerate}
\item 用 \(x_0\)，\(y_0\) 表示点 \(N\) 的坐标 \((x_1, y_1)\)；
\item 用 \(x_0\)，\(y_0\) 和 \(x_1\)，\(y_1\) 表示点 \(P\) 的坐标 \((x, y)\)；
\item 求：当点 \(M\) 沿抛物线移动时，点 \(P\) 的轨迹方程，并道明轨迹是什么图形.
\end{enumerate}
\end{problem}

\begin{solution}
\begin{enumerate}
\item 直线 \(y=\dfrac12x\) 上的点可设为 \(N=(2t,t)\). 因为 \(MN\) 垂直于该直线，而该直线的方向向量为 \((2,1)\)，所以
\[
\bigl(x_0-2t,y_0-t\bigr)\cdot(2,1)=0
\]
解得 \(t=\dfrac{2x_0+y_0}{5}\)，从而 \(x_1=2t=\frac{4x_0+2y_0}{5}\)，\(y_1=t=\frac{2x_0+y_0}{5}\).
\item 因为 \(M\)，\(N\)，\(P\) 依次在同一直线上，且 \(MN:NP=4:1\)，由内分点公式，
\[
N=\frac{M+4P}{5}
\]
因此
\[
P=\frac{5N-M}{4},
\]
即 \(x=\frac{5x_1-x_0}{4}\)，\(y=\frac{5y_1-y_0}{4}\).
\item 将第（1）问的结果代入第（2）问，得 \(x=\frac{3x_0+2y_0}{4}\)，\(y=\frac{x_0}{2}\). 又因为 \(M\) 在抛物线上，所以 \(y_0=x_0^2+2\). 于是 \(x=\frac{2x_0^2+3x_0+4}{4}\)，\(x_0=2y\).
消去 \(x_0\)，得到
\[
x=2y^2+\frac32y+1=2\left(y+\frac38\right)^2+\frac{23}{32}.
\]
由于 \(x_0\) 可以取任意实数，\(y=\dfrac{x_0}{2}\) 也可以取任意实数，所以轨迹是完整的抛物线
\[
x=2\left(y+\frac38\right)^2+\frac{23}{32}.
\]
它的顶点为 \(\left(\dfrac{23}{32},-\dfrac38\right)\)，对称轴为 \(y=-\dfrac38\)，开口方向为 \(x\) 轴正向.
\end{enumerate}
\end{solution}
